\clearpage
\begingroup
\fancyhead[L]{\small THEORIE DER BEWEGUNG DES MONDES.}
\renewcommand{\headrulewidth}{0pt}
\par\noindent\hspace*{\parindent}Wir wollen nunmehro sehen, wie genau dieser Werth der Bewegung des
Knoten mit den Beobachtungen übereinstimme. Wir bemerken daher, dass
\begingroup\small
\setlength{\arraycolsep}{2pt}
\renewcommand{\arraystretch}{1.15}
\[
\begin{array}{rcl@{\qquad}rcl}
\varepsilon&=&3463\overset{\prime\prime}{,}290
 &\log\varepsilon&=&3{,}539\,4888\\
\multicolumn{6}{r}{\text{in Sekunden}}\\
\varepsilon&=&0{,}016\,7905
 &\log\varepsilon&=&8{,}255\,0637\\
\multicolumn{6}{r}{\text{in Theilen des Halbmessers}}\\
e&=&11350\overset{\prime\prime}{,}425
 &\log e&=&4{,}055\,0122\\
\multicolumn{6}{r}{\text{in Sekunden}}\\
e&=&0{,}055\,02843
 &\log e&=&8{,}740\,5871\\
\multicolumn{6}{r}{\text{in Theilen des Halbmessers}}\\
g&=&18583\overset{\prime\prime}{,}18
 &\log g&=&4{,}269\,1200\\
\multicolumn{6}{r}{\text{in Sekunden}}\\
g&=&0{,}090\,09380
 &\log g&=&8{,}954\,6949\\
\multicolumn{6}{r}{\text{in Theilen des Halbmessers}}\\
n&=&0{,}074\,8013\,1353
 &\log n&=&8{,}873\,9092,
\end{array}
\]
\[
\begin{array}[t]{rcl}
1+2ee&=&1{,}006\,0562\,55\\
\tfrac32\varepsilon\varepsilon&=&0{,}000\,4228\,81\\
\hline
&&1{,}006\,4791\,36\\
&&0{,}048\,0203\,88\\
&&0{,}958\,4587\,48
\end{array}
\qquad
\begin{array}[t]{rcl}
\tfrac38n&=&0{,}028\,0504\,92\\
\tfrac12gg&=&0{,}004\,0584\,46\\
\tfrac{91}{32}nn&=&0{,}015\,9114\,50\\
\hline
&&0{,}048\,0203\,88
\end{array}
\]
\endgroup
\SourceLeaf{635}\VIIInlineSourceHead{}{THEORIE DER BEWEGUNG DES MONDES.}{629}
\begingroup\small
\setlength{\arraycolsep}{2pt}
\renewcommand{\arraystretch}{1.15}
\[
\begin{array}{rcll}
\log&=&9{,}981\,5734&\\
\log nn&=&7{,}747\,8184&\\
\log\tfrac34&=&9{,}875\,0613&\\
\cline{1-3}
&&7{,}604\,4531&\\
\log(1+\mu)&=&0{,}000\,0012&\\
\cline{1-3}
&&7{,}604\,4519&=\log0{,}004\,0220\,91.\\
\multicolumn{3}{r}{\text{Die Beobachtungen geben}}&0{,}004\,0218\,746.
\end{array}
\]
\endgroup
\noindent Also der Fehler nur \(\tfrac1{19000}\) des Ganzen.
\par\noindent\hspace*{\parindent}Die Bewegung des Knoten in einem siderischen Jahre findet man, wenn
man diese Grösse mit \(\dfrac{360^\circ}{n}=\dfrac{1296000^{\prime\prime}}{n}\) multiplicirt
\begingroup\small
\setlength{\arraycolsep}{2pt}
\renewcommand{\arraystretch}{1.15}
\[
\begin{array}{rll}
&7{,}604\,4519&\\
\text{Compl. }\log n&1{,}126\,0908&\\
&6{,}112\,6050&\\
\cline{2-2}
&4{,}843\,1477&=\log69686\overset{\prime\prime}{,}35.\\
\multicolumn{2}{r}{\text{Die Beobachtungen geben}}&69682{,}65,
\end{array}
\]
\endgroup
\noindent also der Unterschied \(3\overset{\prime\prime}{,}7\); eine so genaue Übereinstimmung hätte man in der
That kaum erwarten dürfen.
\par\noindent\hspace*{\parindent}Die Integration der Gleichung gibt nunmehro folgenden Werth von \(\theta\):
\[
\theta=g\times
\]
\begingroup\small
\setlength{\arraycolsep}{4pt}
\renewcommand{\arraystretch}{1.3}
\[
\begin{array}{r|l}
\multicolumn{1}{c}{\sin}&\\
\lambda&+1\\
3\lambda&+\tfrac3{32}nngg\\
2E-\lambda&+\tfrac38n(1+\tfrac14n-\tfrac{91}{64}nn+2ee-\tfrac52\varepsilon\varepsilon)\\
2E+\lambda&-\tfrac{75}{128}n^4-\tfrac3{32}nngg\\
2E-3\lambda&-\tfrac3{16}ngg\\
4E-\lambda&+\tfrac{75}{128}n^4\\
M+\lambda&+nne-\tfrac38n^3e\\
M-\lambda&+3nne+\tfrac92n^3e\\
2M+\lambda&+\tfrac{15}{32}nnee
\end{array}
\]
\[
\begin{array}{r|l}
\multicolumn{1}{c}{\sin}&\\
2M-\lambda&+\tfrac58ee\\
2E+M-\lambda&-nne-\tfrac{31}{24}n^3e\\
2E-M+\lambda&+3nne+\tfrac{15}{2}n^3e\\
2E+2M-\lambda&-\tfrac{15}{32}nnee\\
2E-2M-\lambda&-\tfrac{15}{16}nee\\
2E-M+\lambda&+\tfrac{15}{8}n^3e\\
4E-M-\lambda&-\tfrac{15}{8}n^3e\\
\lambda+a&-\tfrac98n\varepsilon(1-\tfrac38n)\\
\lambda-a&+\tfrac98n\varepsilon(1-\tfrac38n)
\end{array}
\]
\endgroup
\SourceLeaf{636}\VIIInlineSourceHead{630}{NACHLASS.}{}
\begingroup\small
\setlength{\arraycolsep}{4pt}
\renewcommand{\arraystretch}{1.3}
\[
\begin{array}{r|l}
\multicolumn{1}{c}{\sin}&\\
\lambda+2a&+\tfrac{27}{32}n\varepsilon\varepsilon\\
\lambda-2a&-\tfrac{27}{32}n\varepsilon\varepsilon\\
2E-\lambda+a&+\tfrac38n\varepsilon(1+\tfrac18n)\\
2E-\lambda-a&-\tfrac78n\varepsilon(1+\tfrac{47}{56}n)\\
2E-\lambda+2a,&*\\
2E-\lambda-2a&+\tfrac{51}{32}n\varepsilon\varepsilon\\
2E-\lambda+M+a&-\tfrac12nne\varepsilon
\end{array}
\]
\[
\begin{array}{r|l}
\multicolumn{1}{c}{\sin}&\\
2E-\lambda+M-a&+\tfrac72nne\varepsilon\\
2E-\lambda-M+a&+\tfrac32nne\varepsilon\\
2E-\lambda-M-a&-\tfrac{21}{2}nne\varepsilon\\
M+\lambda+a&-\tfrac32nne\varepsilon\\
M+\lambda-a&-\tfrac32nne\varepsilon\\
M-\lambda+a&-\tfrac92nne\varepsilon\\
M-\lambda-a&-\tfrac92nne\varepsilon.
\end{array}
\]
\endgroup
\clearpage
\endgroup
